El presente proyecto aborda el desafío logístico de una empresa prestadora de servicios de salud en Santiago de Chile, centrada en la planificación diaria del transporte de más de 960 profesionales hacia y desde sus centros médicos. El problema se define técnicamente como un \textit{Time-Dependent Heterogeneous Dial-a-Ride Problem} (TD-HDARP), cuya complejidad radica en la integración de ventanas de tiempo estrictas, una flota de vehículos variada, y restricciones sindicales de compatibilidad entre pasajeros.

La metodología propuesta consiste en una arquitectura de solución híbrida: un modelo de \textit{Machine Learning (XGBoost)} para capturar la variabilidad del tráfico urbano, y la metaheurística \textit{Adaptive Large Neighborhood Search (ALNS)} para la optimización de rutas. El resultado final será un Sistema de Soporte a la Toma de Decisiones (DSS) desarrollado en Python, que automatiza la planificación y mejora indicadores clave como costos operacionales, ocupación de flota y puntualidad, en un máximo de una hora. El impacto del proyecto trasciende la eficiencia económica, asegurando la continuidad asistencial de los recintos de salud y mejorando el bienestar laboral de los profesionales al reducir su carga logística y mental.